\Qhead{1.7}{Параметрическое усиление и генерация. Деление частоты}{Параметрическое усиление и параметрическая генерация. Параметрические усилители и генераторы. Деление частоты.}
\Subfile{1.7.1}{sections/01_oscillations/q07/s01}
\Subfile{1.7.2}{sections/01_oscillations/q07/s02}

\recap{Накачка модулирует реактивность. Ниже порога прибор усиливает, выше генерирует. Вырожденный усилитель фазочувствителен; невырожденный связывает сигнал с холостой модой. Синхронная генерация на половине накачки даёт делитель на два.}
\board{Варикапный контур; $\omega_p=\omega_s+\omega_i$; квадратуры и порог; потоки энергии Мэнли--Роу.}
\oral{Почему усиление реактивным элементом не нарушает сохранения энергии?}
\bait{Можно оставить два фазовых состояния делителя: обе фазы дают одно и то же удвоенное колебание.}
\sources{Балакший~\cite{balakshy}, PDF-с.~8--9, листы~6--7: параметрическая ёмкость, сигнальный/холостой контуры и схема; Белокопытов~\cite{belokopytov2}, §4.3, PDF-с.~78--80: трёхчастотное модальное взаимодействие. Пятаков~\cite{pyatakov}, PDF-с.~47--52: одноконтурный усилитель. Вопрос~7 коллеги --- основа подачи; коэффициенты и порог получены из выписанной модели.}
